%% kapa: exam=1 round=33 year=2022 period=1 form=A subject=real_estate subject_ko=부동산학원론 seq=17 number=97
%% source: Q-Net 감정평가사 기출문제 (cst00309) artlSeq=5209997 file=raw/1008/5209997/2022년도 제33회 감정평가사 1차시험 1교시(민법,경제학,부동산학).pdf
%% conversion: kordoc 4.24.51 -> tools/gen1.py (auto); review=proofread-sonnet; flags=none
\begin{question}{97}
대출상환방식에 관한 설명으로 옳지 않은 것은? (단, 주어진 조건에 한함)
\begin{choices}
  \item 원금균등분할상환방식은 만기에 가까워질수록 차입자의 원리금상환액이 감소한다.
  \item 원리금균등분할상환방식은 만기에 가까워질수록 원리금상환액 중 원금의 비율이 높아진다.
  \item 대출조건이 동일하다면 대출기간동안 차입자의 총원리금상환액은 원금균등분할상환방식이 원리금균등분할상환방식보다 크다.
  \item 차입자의 소득에 변동이 없는 경우 원금균등상환방식의 총부채상환비율(DTI)은 만기에 가까워질수록 낮아진다.
  \item 차입자의 소득에 변동이 없는 경우 원리금균등분할상환방식의 총부채상환비율은 대출기간동안 일정하게 유지된다.
\end{choices}
\end{question}
\answer{3}
